\subsection{Controlled Room study}\label{sec:room-study}
The follow-up replays 7,000 fitting updates at seed 42 with the current code, the archived fitting camera order and the same 39 evaluation targets. Both primary arms use zero loader workers. All roles share the same fitting objective, initialization and camera order. Every arm is audited against the known initializer and the target RGB hashes. This study is a new measurement. The original five-scene comparisons above retain their own endpoint values, and comparisons in this section use only runs from the same replay.

The Room protocol crosses gradient aggregation (signed or absolute) with eligibility threshold ($2\times10^{-4}$ or $8\times10^{-4}$) under the same cap, ranking and split operator. It also disables RGB-error attribution and edge strength one at a time. Three cap settings measure the realized size and quality at the default absolute threshold. The cap sweep showed that the 700,000 ceiling did not restrict A05-Abs selection on Room. Two additional runs remove each ranking term at the binding 400,000 cap, where the remaining scores decide which eligible parents receive the available slots.

\subsubsection{Matched updates and elapsed training time}
The main pair also evaluates after updates 1, 1,000, 2,000, \ldots, 7,000 and at target elapsed times 90, 180 and 270 seconds. Time starts at the training loop after loading and initialization. CUDA synchronization measures completed work. Deadline checks occur at 10-update intervals, and elapsed time excludes evaluation. Each sample is taken just after crossing its time target. Table~\ref{tab:matched-time} reports the actual times and counts. The same elapsed budget can correspond to a different number of optimizer updates. Figure~\ref{fig:room-progress} plots the trajectories.

\begin{center}\begin{minipage}{\linewidth}
\small\begin{tabular*}{\linewidth}{@{\extracolsep{\fill}}rlrrrrrr@{}}
\toprule
Target (s) & Arm & Actual (s) & Updates & PSNR (dB) & SSIM & LPIPS & Gaussians\\
\midrule
\inputrows{evidence/room_time_rows.tex}
\bottomrule
\end{tabular*}\par
\captionof{table}[Room results at matched elapsed training times]{Public Room at shared elapsed training targets. Time excludes interleaved evaluation.}\label{tab:matched-time}
\end{minipage}\end{center}

\reportfigure[Room learning curves by update count and elapsed time]{0.98\linewidth}{figures/room_progress.pdf}
  {Room learning trajectories. The update panel uses fixed update milestones. The time panel uses the actual seconds measured near the three common deadlines. Both evaluate the same 39 development views. The update panel omits the nearly identical first-update scores to show the later differences.}{fig:room-progress}
At the 90, 180 and 270 second targets, the candidate led by 0.099, 0.261 and 0.513 dB, respectively, with fewer Gaussians at each sample. At 90 seconds its SSIM and LPIPS were slightly worse. At 270 seconds it completed 5,290 updates versus 5,390 for B00, while using 432,794 versus 781,109 Gaussians and scoring 30.236 versus 29.723 dB. The candidate accumulated splats more quickly through the first 1,000 updates. Between 1,000 and 2,000 updates, the baseline count overtook it, while the candidate maintained a PSNR lead. At the sampled deadlines, the PSNR gain and count reduction occurred despite fewer completed updates.

The final replay scores were 30.757 dB at 452,409 Gaussians for A05-Abs and 30.125 dB at 834,493 for B00. This is a 0.632 dB difference and a 45.8\% count reduction. Both scores continued to rise between the 6,000- and 7,000-update probes with counts unchanged, reflecting attribute optimization after structural refinement stopped. The paired replay reproduced the direction of the archived Room result. Exact scores and counts differed slightly from the archive, so the replay and archived values are reported separately [\hyperlink{evE8}{E8}].


\subsubsection{Gradient aggregation, threshold and ranking}
For a metric $M$, let $M_{a,h}$ denote the absolute-gradient result at the higher threshold, with $s$ denoting signed and $l$ the lower threshold. The effect of changing aggregation at fixed high threshold is $M_{a,h}-M_{s,h}$. The threshold effect within absolute gradients is $M_{a,h}-M_{a,l}$. The interaction is $(M_{a,h}-M_{s,h})-(M_{a,l}-M_{s,l})$. Table~\ref{tab:room-components} gives the four single-seed arms. Table~\ref{tab:tight-ranking} repeats the ranking contrasts at a binding ceiling. Counts and appearance jointly reflect the accumulated changes to eligibility and parent ranking.

\begin{center}\begin{minipage}{\linewidth}
\small\begin{tabular*}{\linewidth}{@{\extracolsep{\fill}}llrrrrr@{}}
\toprule
Policy & Aggregation & Threshold & PSNR (dB) & SSIM & LPIPS & Gaussians\\
\midrule
\inputrows{evidence/room_component_rows.tex}
\bottomrule
\end{tabular*}\par
\captionof{table}[Room gradient and ranking ablations at a 700,000 cap]{Room gradient-by-threshold factorial and image-ranking ablations at a 700,000-Gaussian cap. All rows use Long-Axis Split and 7,000 updates.}\label{tab:room-components}
\end{minipage}\end{center}
\begin{center}\begin{minipage}{\linewidth}
\small\begin{tabular*}{\linewidth}{@{\extracolsep{\fill}}lrrrrr@{}}
\toprule
Ranking policy & Binding events & PSNR (dB) & SSIM & LPIPS & Gaussians\\
\midrule
\inputrows{evidence/room_ranking_rows.tex}
\bottomrule
\end{tabular*}\par
\captionof{table}[Ranking ablations at a binding 400,000 cap]{Ranking ablations at the binding 400,000-Gaussian cap. All arms use absolute gradients and threshold $8\times10^{-4}$.}\label{tab:tight-ranking}
\end{minipage}\end{center}
Absolute aggregation raised PSNR by 0.487 dB at the low threshold and by 1.089 dB at the high threshold, an interaction of about 0.602 dB. It also increased count at each fixed threshold. The absolute low-threshold arm reached 31.084 dB and 699,645 Gaussians, with competition for slots at all 54 events. The selected high-threshold arm gave up 0.327 dB and saved 247,236 primitives (35.3\%). It still improved on B00's 30.125 dB by 0.632 dB with 45.8\% fewer primitives. The high-threshold signed arm grew to only 143,086 Gaussians and fell below B00 on all three RGB metrics. These contrasts measure the gradient eligibility policy under fixed split and ranking settings. The low threshold also activates the cap, so the threshold contrast includes that downstream effect.

The earlier A04 factorial supports a contribution from Long-Axis Split and replacing clones with splits under signed gradients (Table~\ref{tab:split-ablation}). An equivalent factorial under A05-Abs would measure whether the split effects change with absolute gradients [\hyperlink{evE3}{E3}].

At 700,000, every eligible parent split in all three absolute/high-threshold ranking arms. At 400,000, competition occurred in 32 or 33 events. Removing the RGB-error cue changed PSNR by $-0.018$ dB and removing the edge cue changed it by $+0.033$ dB, at nearly equal counts. Both removals slightly increased LPIPS. The full configuration remains the fixed reference because it is the design carried forward from development and has the lowest LPIPS in this binding-cap trio. The edge-free variant has the highest PSNR here. A single Room seed and opposing metric directions leave the general value of the cues unresolved [\hyperlink{evE8}{E8}].


\subsubsection{Gaussian cap and realized count}\label{sec:room-caps}
Sweeping the cap tests how the same strategy behaves under different count ceilings. Each setting uses the same gradient definition and threshold, ranking, split operation and 7,000 training updates. The cap changes the event targets throughout training. It can alter which parents survive and compete later, so two similar endpoints can still arise through different trajectories. The baseline contributes one unrestricted endpoint for context.

\begin{center}\begin{minipage}{\linewidth}
\small\begin{tabular*}{\linewidth}{@{\extracolsep{\fill}}lrrrrr@{}}
\toprule
Arm & Cap & Gaussians & PSNR (dB) & SSIM & LPIPS\\
\midrule
\inputrows{evidence/room_cap_rows.tex}
\bottomrule
\end{tabular*}\par
\captionof{table}[Room cap sweep at equal update count]{Room cap sweep at equal update count. The final count is the measured model size.}\label{tab:room-caps}
\end{minipage}\end{center}

\reportfigure[Room PSNR against realized Gaussian count]{0.92\linewidth}{figures/room_cap_tradeoff.pdf}
  {Room quality against realized Gaussian count at 7,000 updates. The three budgeted points highlight the tighter cap, the selected high threshold and the low threshold. Table~\ref{tab:room-caps} also gives the 1,000,000-ceiling control, which overlaps the selected operating point at this scale.}{fig:cap-tradeoff}
At the selected threshold, both the 700,000 and 1,000,000 schedules admitted every eligible parent and finished near 450,000 Gaussians. Top-$k$ ranking and cap enforcement excluded zero parents in the selected 700,000 run. Its gain over B00 measures the bundled changes in eligibility and splitting. The exact contribution of each change remains unresolved. The ceiling limits growth where eligibility exceeds capacity.

At 400,000, the guardrail constrained 33 of 54 events. It saved 52,590 primitives relative to the 700,000-cap arm at a cost of 0.051 dB. It still led B00 by 0.581 dB with 52.1\% fewer primitives. Figure~\ref{fig:cap-tradeoff} shows the resulting quality--representation-size choices. A larger model bought higher PSNR at the low absolute threshold, and a stricter cap bought a smaller model at a measurable quality cost [\hyperlink{evE8}{E8}].

